% ============================================================
\chapter{Programmation en Nombres Entiers}
\label{ch:nombres_entiers}
% ============================================================

\section{Introduction}

En programmation lin\'eaire, les solutions optimales vivent dans un monde continu~: on peut affecter $3{,}7$ camions \`a une route. Mais dans le monde r\'eel, un camion est entier. On ne peut pas couper un avion en deux pour desservir deux lignes. Cette contrainte d'int\'egralit\'e, apparemment anodine, transforme radicalement la nature du probl\`eme~: la programmation en nombres entiers est NP-difficile dans le cas g\'en\'eral, et les m\'ethodes de r\'esolution sont fondamentalement diff\'erentes. La m\'ethode de \emph{branch and bound}, d\'evelopp\'ee par Ailsa Land et Alison Doig en 1960, explore intelligemment un arbre de sous-probl\`emes~; les \emph{coupes de Gomory} (1958) ajoutent des contraintes pour resserrer la relaxation continue~; le \emph{branch and cut} combine les deux. Ces algorithmes sont au c\oe ur des solveurs modernes (CPLEX, Gurobi) qui r\'esolvent des probl\`emes industriels \`a des millions de variables.

\begin{intuition}
La programmation lin\'eaire continue produit des solutions qui peuvent \^etre
fractionnaires (par exemple, fabriquer 3.7 tables). Arrondir na\"ivement la
solution n'est en g\'en\'eral ni optimal ni m\^eme r\'ealisable. La PNE
fournit des m\'ethodes syst\'ematiques pour trouver la meilleure solution
enti\`ere.
\end{intuition}

\section{Formulation}

\begin{definition}[Programme lin\'eaire en nombres entiers]
\index{programmation en nombres entiers}
Un \textbf{programme lin\'eaire en nombres entiers} (PLNE) est~:
\[
  \min\; c^\top x \quad \text{s.c.} \quad Ax \leq b, \quad x \in \Z^n_+.
\]
Si seules certaines variables sont enti\`eres, on parle de \textbf{programme
mixte} (PLNEM). Si toutes les variables sont binaires ($x \in \{0,1\}^n$), on
parle de \textbf{programme en 0-1}.
\end{definition}

\begin{definition}[Relaxation lin\'eaire]
\index{relaxation lin\'eaire}
La \textbf{relaxation lin\'eaire} d'un PLNE est le programme lin\'eaire obtenu
en rempla\c{c}ant $x \in \Z^n_+$ par $x \in \R^n_+$. Sa valeur optimale
fournit une \textbf{borne inf\'erieure} (pour un probl\`eme de minimisation)
de la valeur optimale du PLNE.
\end{definition}

\begin{proposition}[Borne de relaxation]
Soit $z^*_{\mathrm{LP}}$ la valeur optimale de la relaxation et
$z^*_{\mathrm{IP}}$ celle du PLNE. Alors~:
\[
  z^*_{\mathrm{LP}} \leq z^*_{\mathrm{IP}}.
\]
Si la solution optimale de la relaxation est enti\`ere, elle est aussi optimale
pour le PLNE.
\end{proposition}

\begin{formules}[title=\textbf{Mod\'elisation classique en 0-1}]
\begin{itemize}
  \item \textbf{S\'election}~: $x_j = 1$ si l'objet $j$ est s\'electionn\'e.
  \item \textbf{Implication}~: $x_i = 1 \Rightarrow x_j = 1$ se mod\'elise
        par $x_i \leq x_j$.
  \item \textbf{Disjonction}~: au moins un parmi $x_1, \ldots, x_k$~:
        $x_1 + \cdots + x_k \geq 1$.
  \item \textbf{Big-M}~: $y \leq M x$ force $y = 0$ quand $x = 0$.
\end{itemize}
\end{formules}

\section{M\'ethode de Branch and Bound}

\begin{definition}[Branch and Bound]
\index{branch and bound}
La m\'ethode de \textbf{Branch and Bound} (s\'eparation et \'evaluation)
r\'esout un PLNE par \'enum\'eration implicite en explorant un arbre de
recherche. \`A chaque n\oe ud, on r\'esout la relaxation lin\'eaire et on
utilise des r\`egles d'\'elagage pour \'eviter d'explorer des branches
sous-optimales.
\end{definition}

\begin{algorithme}[title=\textbf{Branch and Bound}]
\begin{enumerate}
  \item R\'esoudre la relaxation lin\'eaire du probl\`eme initial.
  \item Si la solution est enti\`ere, c'est l'optimum. \textbf{Stop}.
  \item Sinon, choisir une variable $x_j$ fractionnaire (valeur $\bar{x}_j$).
  \item \textbf{S\'eparation}~: cr\'eer deux sous-probl\`emes~:
    \begin{itemize}
      \item Branche gauche~: ajouter $x_j \leq \lfloor \bar{x}_j \rfloor$.
      \item Branche droite~: ajouter $x_j \geq \lceil \bar{x}_j \rceil$.
    \end{itemize}
  \item \textbf{\'Evaluation}~: r\'esoudre la relaxation de chaque
        sous-probl\`eme.
  \item \textbf{\'Elagage}~: \'eliminer un n\oe ud si~:
    \begin{itemize}
      \item le sous-probl\`eme est infaisable,
      \item sa borne inf\'erieure $\geq$ meilleure solution connue,
      \item sa solution est enti\`ere (mettre \`a jour la meilleure solution).
    \end{itemize}
  \item R\'ep\'eter jusqu'\`a ce que tous les n\oe uds soient \'elagu\'es.
\end{enumerate}
\end{algorithme}

\begin{exemple}[Branch and Bound]
Consid\'erons le PLNE~:
\[
  \max\; 5x_1 + 4x_2 \quad \text{s.c.} \quad
  6x_1 + 4x_2 \leq 24, \quad x_1 + 2x_2 \leq 6, \quad
  x_1, x_2 \in \Z_+.
\]
La relaxation donne $x_1 = 3.6$, $x_2 = 0.6$, $z = 20.4$.

On s\'epare sur $x_1$~:
\begin{itemize}
  \item Branche $x_1 \leq 3$~: optimum LP en $x_1 = 3, x_2 = 1.5$,
        $z = 21$.
  \item Branche $x_1 \geq 4$~: optimum LP en $x_1 = 4, x_2 = 0$,
        $z = 20$ (solution enti\`ere).
\end{itemize}
On continue sur la branche gauche en s\'eparant sur $x_2$~:
$x_2 \leq 1$ donne $x_1 = 3, x_2 = 1$, $z = 19$ (enti\`ere) et
$x_2 \geq 2$ donne $x_1 = 2\tfrac{2}{3}, x_2 = 2$, $z \approx 21.3$.
On s\'epare encore~: $x_1 \leq 2$ donne $z = 18$~; $x_1 \geq 3$ est
infaisable. L'optimum est $x_1 = 4, x_2 = 0$ avec $z^* = 20$.
\end{exemple}

\section{M\'ethodes de coupes}

\begin{definition}[Coupe valide]
\index{coupe valide}
Une \textbf{coupe valide} (ou in\'egalit\'e valide) est une contrainte
lin\'eaire satisfaite par toutes les solutions enti\`eres r\'ealisables mais
qui \'elimine la solution fractionnaire courante de la relaxation.
\end{definition}

\subsection{Coupes de Gomory}

\begin{theoreme}[Coupe de Gomory]
\index{coupes de Gomory}
Soit $x_i = \bar{b}_i - \sum_j \bar{a}_{ij} x_j$ une ligne du tableau
simplexe final (avec $\bar{b}_i \notin \Z$). Alors la coupe~:
\[
  \sum_j \{\bar{a}_{ij}\} \, x_j \geq \{\bar{b}_i\}
\]
est valide, o\`u $\{y\} = y - \lfloor y \rfloor$ d\'esigne la partie
fractionnaire.
\end{theoreme}

\begin{proof}
La ligne du tableau donne $x_i + \sum_j \bar{a}_{ij} x_j = \bar{b}_i$.
Comme $x_i \in \Z$ et les variables hors base $x_j \geq 0$, on d\'ecompose~:
\[
  \underbrace{x_i + \sum_j \lfloor \bar{a}_{ij} \rfloor x_j
  - \lfloor \bar{b}_i \rfloor}_{\in \Z}
  = \{\bar{b}_i\} - \sum_j \{\bar{a}_{ij}\} x_j.
\]
Le membre gauche est entier. Or $\{\bar{b}_i\} > 0$ et chaque
$\{\bar{a}_{ij}\} x_j \geq 0$. Pour que le membre droit soit un entier
non-n\'egatif (\`a droite, le terme $x_i + \cdots$ pourrait \^etre n\'egatif),
on obtient $\sum_j \{\bar{a}_{ij}\} x_j \geq \{\bar{b}_i\}$ en
consid\'erant le cas le plus restrictif.
\end{proof}

\begin{remarque}
L'algorithme de Gomory ajoute it\'erativement des coupes \`a la relaxation
lin\'eaire. En th\'eorie, un nombre fini de coupes suffit pour les PLNE purs.
En pratique, les m\'ethodes de \textbf{Branch and Cut} combinent Branch and
Bound et g\'en\'eration de coupes \`a chaque n\oe ud.
\end{remarque}

\section{Unimodularit\'e totale}

\begin{definition}[Matrice totalement unimodulaire]
\index{unimodularit\'e totale}
Une matrice $A \in \Z^{m \times n}$ est \textbf{totalement unimodulaire} (TU)
si le d\'eterminant de toute sous-matrice carr\'ee est dans $\{-1, 0, 1\}$.
\end{definition}

\begin{theoreme}[Relaxation exacte]
Si $A$ est totalement unimodulaire et $b \in \Z^m$, alors le poly\`edre
$\{x \in \R^n : Ax \leq b, x \geq 0\}$ a tous ses sommets entiers. La
relaxation lin\'eaire fournit donc automatiquement une solution enti\`ere.
\end{theoreme}

\begin{proof}
Tout sommet du poly\`edre est solution d'un syst\`eme $A' x = b'$ o\`u
$A'$ est une sous-matrice carr\'ee inversible de la matrice des contraintes
actives. Comme $A$ est TU, $\det(A') \in \{-1, 1\}$ et par la r\`egle de
Cramer, $x = (A')^{-1} b'$ est entier puisque
$(A')^{-1} = \frac{1}{\det(A')} \mathrm{adj}(A')$ avec $\mathrm{adj}(A')$
enti\`ere.
\end{proof}

\begin{exemple}[Matrices TU classiques]
\begin{itemize}
  \item La matrice d'incidence sommets-arcs d'un graphe orient\'e est TU.
  \item La matrice de contraintes du probl\`eme de transport est TU.
  \item Cons\'equence~: les probl\`emes de flots et de transport se r\'esolvent
        par le simplexe et donnent des solutions enti\`eres.
\end{itemize}
\end{exemple}

\section{Applications classiques}

\subsection{Probl\`eme du sac \`a dos}

\begin{definition}[Sac \`a dos 0-1]
\index{sac \`a dos}
\'Etant donn\'es $n$ objets de valeurs $v_1, \ldots, v_n$ et de poids
$w_1, \ldots, w_n$, et un sac de capacit\'e $W$~:
\[
  \max \sum_{j=1}^n v_j x_j \quad \text{s.c.} \quad
  \sum_{j=1}^n w_j x_j \leq W, \quad x_j \in \{0, 1\}.
\]
\end{definition}

\subsection{Probl\`eme du voyageur de commerce}

\begin{definition}[TSP -- formulation]
\index{voyageur de commerce (TSP)}
Le \textbf{probl\`eme du voyageur de commerce} (TSP) sur $n$ villes~:
\[
  \min \sum_{i \ne j} d_{ij} x_{ij} \quad \text{s.c.} \quad
  \sum_{j \ne i} x_{ij} = 1 \;\forall i, \quad
  \sum_{i \ne j} x_{ij} = 1 \;\forall j, \quad
  x_{ij} \in \{0,1\},
\]
plus des contraintes d'\'elimination de sous-tours.
\end{definition}

\begin{attention}[title=\textbf{Complexit\'e du TSP}]
Le TSP est NP-difficile. Les solveurs modernes (Concorde) utilisent des
m\'ethodes de Branch and Cut avanc\'ees et peuvent r\'esoudre des instances
de plusieurs milliers de villes.
\end{attention}

\section{Exercices}

\begin{exercice}[$\star$ -- Branch and Bound]
R\'esoudre par Branch and Bound~:
$\max\; 3x_1 + 2x_2$ s.c. $2x_1 + x_2 \leq 6$, $x_1 + 3x_2 \leq 9$,
$x_1, x_2 \in \Z_+$. Dessiner l'arbre d'exploration.
\end{exercice}

\begin{exercice}[$\star\star$ -- Coupes de Gomory]
Consid\'erer $\max\; x_1 + x_2$ s.c. $3x_1 + 2x_2 \leq 6$,
$x_1 + 4x_2 \leq 4$, $x_1, x_2 \geq 0$ entiers. G\'en\'erer une coupe de
Gomory \`a partir du tableau simplexe final et r\'esoudre.
\end{exercice}

\begin{exercice}[$\star\star$ -- Sac \`a dos]
R\'esoudre le sac \`a dos 0-1~: objets de valeurs $(10, 6, 12, 7)$, poids
$(3, 2, 5, 3)$, capacit\'e $W = 8$. Comparer Branch and Bound et
programmation dynamique.
\end{exercice}

\begin{exercice}[$\star\star\star$ -- Formulation TSP]
Formuler le TSP sur 5 villes avec les contraintes de Miller-Tucker-Zemlin
(MTZ). R\'esoudre la relaxation lin\'eaire et mesurer le \emph{gap}
d'int\'egralit\'e.
\end{exercice}

\begin{exercice}[$\star\star\star$ -- Unimodularit\'e]
Montrer que la matrice d'incidence sommets-arcs d'un graphe orient\'e est
totalement unimodulaire. En d\'eduire que le probl\`eme de flot \`a co\^ut
minimum admet toujours une solution enti\`ere optimale quand les donn\'ees
sont enti\`eres.
\end{exercice}
